% Terjemahan dibuat oleh AI. Koreksi dan tata letak dihasilkan oleh OpenAI Codex — GPT-6 Astra, tingkat penalaran Ultra. Nama Emmy Noether dan para penulis yang tercantum pada karya tetap dipertahankan. Tidak ada penerjemahan manusia atau penelaahan ilmiah manusia independen yang diatribusikan kepada penerbit.
% Hak tetap dimiliki pemegangnya masing-masing. Penyebaran gratis untuk tujuan nonkomersial dan murni pendidikan didasarkan pada anggapan penerbit bahwa pemegang hak tidak keberatan, bukan izin yang telah dipastikan. Akses gratis tidak dengan sendirinya mengizinkan publikasi turunan. Pemegang hak dapat mengajukan keberatan melalui tautan berikut; penerbit akan menarik materi terkait atas permintaan mereka. Tidak ada lisensi baru atau penyerahan ke domain publik.
% https://github.com/KokunoYumeto/modern-latex-manuscripts/issues/new
\documentclass[11pt,a4paper]{article}
\usepackage[margin=28mm]{geometry}
\usepackage{fontspec}
\setmainfont{TeX Gyre Pagella}
\usepackage[indonesian]{babel}
\usepackage{microtype}
\setlength{\parindent}{1.4em}
\setlength{\parskip}{0pt}
\begin{document}
\begin{center}
{\Large\bfseries 6. Lapangan dan Sistem Fungsi Rasional}\par
\vspace{0.6em}
Emmy Noether (Erlangen).\par
\vspace{1.2em}
\emph{Math. Ann.} 76 (1915), hlm. 161--196
\end{center}

\vspace{2em}

Tulisan ini membahas persoalan-persoalan basis pada sistem sebarang fungsi rasional dan fungsi rasional integral. Metode teori lapangan yang digunakan memperlihatkan bahwa penyelesaian persoalan-persoalan ini untuk lapangan yang terdiri atas fungsi rasional---yakni lapangan fungsi rasional\footnote{Bandingkan pengumuman pendahuluan saya pada pertemuan ilmuwan alam di Wina: ``Rationale Funktionenkörper'', \emph{Jahresber. d. D. Math.-Ver.} 22 (1913), yang memberikan tinjauan atas persoalan-persoalan tersebut dan hasil-hasil mengenai lapangan fungsi.}---merupakan bagian yang hakiki, sedangkan perluasan hasil-hasil itu ke sistem sebarang diperoleh sebagai konsekuensi.

Mengenai persoalan basis pada sistem umum, sejauh ini hanya eksistensi basis modul yang dijamin oleh teorema Hilbert (\emph{Math. Ann.} 36) yang diketahui. Dalam uraian berikut, persoalan keterwakilan secara rasional dijawab sepenuhnya melalui eksistensi basis rasional bagi setiap sistem sebarang (\S\ 7); basis rasional untuk lapangan telah diperoleh dalam \S\ 4. Eksistensi basis rasional ini memungkinkan kita untuk selalu bertolak dari lapangan atau sistem yang didefinisikan secara abstrak, dan dengan demikian menghindari kesulitan yang semata-mata timbul dari pemilihan khusus basis rasional, bukan dari sistem itu sendiri, misalnya munculnya penyebut khusus atau titik fundamental fungsi-fungsi basis.
\end{document}
